\specsection{Brown 运动的理论}

这一理论在当代科学发展中占有特殊的地位，它见证了物理问题与“纯”数学之间持续而富有成果的交流。Brown 运动由植物学家 Brown 于 1829 年发现，19 世纪众多物理学家对它进行了深入的研究，\(^{(3)}\) 但第一个令人满意的数学模型直到 1905 年才由 Einstein 发明。在粒子沿直线运动的简单情形中，Einstein 的基本假设可以表述如下：若 \(x(t)\) 是粒子在时刻 t 的横坐标，且 \(t_{0} < t_{1} < \cdots < t_{n-1} < t_{n}\)，则相继的位移 \(x(t_{i}) - x(t_{i-1})\)（对 \(1 \leqslant i \leqslant n\)）是独立的 Gauss 随机变量。这里不是详细评述促使 Einstein 理论产生的 J. Perrin 的重要实验工作的地方；对我们的目的而言，只需记住 Perrin 的一个注记，按照这个注记，Brown 运动轨迹的观察使他不可抗拒地联想到“数学家们的没有导数的函数”。这个注记成为 Wiener 的最初火花。

另一个思想潮流发源于气体动理论，它由 Boltzmann 与 Gibbs 在 1870 年至 1900 年间发展。考虑由 N 个质量为 m 的分子组成的气体，处于（绝对）温度 T，并用 \(v_{1},\ldots,v_{N}\) 表示该气体 N 个分子的速度；系统的动能等于

\[
\frac {m}{2} (\mathbf {v} _ {1} ^ {2} + \dots + \mathbf {v} _ {\mathrm{N}} ^ {2}) = 3 \mathrm{N} k \mathrm{T},\tag{1}
\]

其中 k 是 Boltzmann 常数。按照 Gibbs 的想法，分子之间大量的碰撞不容许精确地确定分子的速度，于是方便的做法是在由方程 (1) 定义的 3N 维空间的球面 S 上引入一个概率律 P。“微正则”假设在于假定 P 是球面 S 上在旋转下不变的质量为 1 的测度。此外，Maxwell 速度律指出，分子速度的一个分量的概率律是一个方差为 \(2k\mathrm{T} / m\) 的 Gauss 测度（§6, No.~5, Remark 3）。Borel 似乎是最早在 1914 年观察到：当分子数很大时，Maxwell 律是 Gibbs 假设与球面性质的一个推论。他考虑高维 Euclid 空间中的一个球面 S 以及 S 上在旋转下不变的质量为 1 的测度 P；利用基于 Stirling 公式的经典近似方法，他证明了 P 在一条坐标轴上的投影近似地是 Gauss 的。这些结果稍后由 Gâteaux 与 Lévy (IX, \(a\) ) 加以精确化。给定一个整数 \(m \geqslant 1\) 与一个数 \(r > 0\)，用 \(\mathrm{S}_{m,r}\) 表示形如 \((x_{1}, \ldots, x_{m}, 0, 0, \ldots)\)、满足 \(x_{1}^{2} + \cdots + x_{m}^{2} = r^{2}\) 的序列所成的集合；再用 \(\sigma_{m,r}\) 表示 \(\mathrm{S}_{m,r}\) 上在旋转下不变的质量为 1 的测度。用现代的语言陈述，Gâteaux 与 Lévy 的结果如下：测度序列 \(\sigma_{m,1}\) 紧收敛到原点 \((0, 0, \ldots)\) 处的一个单位质量，而测度序列 \(\sigma_{m,\sqrt{m}}\) 紧收敛到一个形如下式的测度 \(\Gamma\)

\[
d \Gamma (x _ {1}, x _ {2}, \dots) = \prod_ {n = 1} ^ {\infty} d \gamma (x _ {n})
\]

（γ 是 R 上方差为 1 的 Gauss 测度）。

前面的测度 \(\Gamma\) 起着无穷维 Gauss 测度的作用。Lévy 似乎确实曾模糊地希望以内在的方式在每个无穷维 Hilbert 空间上定义一个 Gauss 测度。事实上，正如 Lévy 与 Wiener 所指出的，测度 \(\Gamma\) 在某种意义下 \(^{(4)}\) 在 \(l^{2}\) 的自同构之下是不变的；不幸的是，集合 \(l^{2}\)（由平方可和序列 \((x_{1}, x_{2}, \ldots, x_{n}, \ldots)\) 组成）对 \(\Gamma\) 的测度为零。现在人们知道，只能满足于无穷维 Hilbert 空间上的一个 Gauss 预测度。\(^{(5)}\)

本质的进展要归功于 Wiener：既然在无穷维 Hilbert 空间上没有合理的 Gauss 测度，人们就可以借助取原函数的运算，从一个 Gauss 预测度出发，在连续函数空间上构造一个测度 \(w\)（细节 cf.~§6, No.~7, Th. 1）。我们将扼要地解释 Wiener 对 \(w\) 的原始构造 (X)；它直接受到 Gâteaux 与 Lévy 的关系式 \(\Gamma = \lim_{m \to \infty} \sigma_{m,\sqrt{m}}\) 的影响。对每个整数 \(m \geqslant 1\)，用 \(H_m\) 表示 \(T = ]0,1]\) 上在每个区间 \(\left[ \frac{k-1}{m}, \frac{k}{m} \right]\)（对 \(k = 1,2,\ldots,m\)）上取常值的函数所成的集合，用 \(\pi_m\) 表示 \(R^m\) 中半径为 1 的 Euclid 球面上在旋转下不变的质量为 1 的测度。设 \(f_m\) 是 \(H_m\) 到 \(R^m\) 上的同构，它把取值 \(a_k\)（在区间 \(\left[ \frac{k-1}{m}, \frac{k}{m} \right]\) 上）的函数对应到向量 \((a_1, a_2 - a_1, \ldots, a_m - a_{m-1})\)（Wiener 钟爱的“微分空间”一词即由此而来）；用 \(w_m\) 表示 \(H_m\) 上 \(\pi_m\) 在 \(f_m^{-1}\) 下的像测度。Wiener 把所想要的测度 \(w\) 定义为测度 \(w_m\) 的极限。确切地说，用 \(H\) 表示 \(T\) 上赋以一致收敛拓扑的调节函数所成的集合（有 \(H_m \subset H\)，对每个整数 \(m \geqslant 1\)）；对每个有界的一致连续函数 \(F\)（\(H\) 上的），极限 \(A\{F\} = \lim_{m \to \infty} \int_{H_m} F(x) dw_m(x)\) 存在；接着，Wiener 通过对掷硬币游戏之起伏的精细分析得到某些上界估计，并重新利用 Daniell 所强调的紧性论证，证明此时满足应用 Daniell 扩张定理的条件。由此得出存在一个测度 \(w\)（由 \(C(T)\) 支撑），使得 \(A\{F\} = \int_{C(T)} F(x) dw(x)\)。Wiener 随后能够证明测度 \(w\) 对应于 Einstein 的假设，\(^{(6)}\) 而且他的估计使他能够对 Perrin 关于无导数函数的注记赋予精确的意义：满足阶为 \(\frac{1}{2}\) 的 Lipschitz 条件的函数所成的集合对 \(w\) 是可忽略的（然而，对每个 \(a\)（满足 \(0 < a < \frac{1}{2}\)），几乎每个函数都满足阶为 \(a\) 的 Lipschitz 条件）。

如今，Wiener 测度的众多构造已是熟知的。例如，Paley 与 Wiener 使用了随机 Fourier 级数 (XI, Ch. IX)：对每个实

\[
\begin{array}{l} \int_ {\mathcal {C} (\mathrm{T})} f (x (t _ {1}), \dots , x (t _ {n})) d w (x) = \\ (2 \pi) ^ {- n / 2} \prod_ {i = 1} ^ {n} (t _ {i} - t _ {i - 1}) ^ {- 1 / 2} \int \dots \int f (x _ {1}, \dots , x _ {n}) \exp \left(- \frac {1}{2} \sum_ {i = 1} ^ {n} \frac {(x _ {i} - x _ {i - 1}) ^ {2}}{t _ {i} - t _ {i - 1}}\right) d x _ {1} \dots d x _ {n}, \end{array}
\]

序列 \(\mathbf{a} = (a_n)_{n\geqslant 1}\) 与每个整数 \(m\geqslant 0\)，用下式定义 ]0,1] 上的函数 \(f_{m,\mathbf{a}}\)：

\[
f _ {m, \mathbf {a}} (t) = a _ {1} t + 2 \sum_ {k = 2} ^ {2 ^ {m + 1}} \frac {1}{\pi k} a _ {k - 1} \sin \pi k t;
\]

可以证明，对 \(\Gamma\)-几乎每个序列 \(\mathbf{a}\)，函数序列 \(f_{m,\mathbf{a}}\) 趋于一个连续函数 \(f_{\mathbf{a}}\)，且 \(w\) 是 \(\Gamma\) 在（几乎处处定义的）映射 \(\mathbf{a} \mapsto f_{\mathbf{a}}\) 下的像。后来，Lévy 在 (IX, \(b\) ), \(c\) ) 中给出了一个与我们在 §6, No.~7 中所阐述的非常接近的构造。最后，Kac、Donsker 与 Erdős 大约在 1950 年指明，如何用更一般的测度代替 Wiener 原始构造中的球面测度 \(\pi_m\)（在 \(\mathbf{R}^m\) 上）。他们的结果在 Wiener 测度与概率论的极限定理之间建立了牢固的联系；这些结果后来由 Prokhorov (XIII) 加以完成和系统化，我们将在后面回到这一工作。

这里不是分析由 Wiener 的发现所产生的众多而重要的概率论工作的地方；如今，Brown 运动只表现为 Markoff 过程最重要的一些例子之一。我们只提一下 Kac 把 Wiener 测度应用于求解某些抛物型偏微分方程；这是对量子场论中 Feynmann 思想的一个改编---这又是数学与物理问题相互影响的一个例子。

